\documentclass{zhvt-classic}

% 使用本地源流明体（GenRyuMin2 TW）按文件路径加载，不依赖系统字体名
% 版框（内框／界栏／书口侧栏／鱼尾）改用与版心印章同一印泥棕色：
%   印章图片 yshs.png 实色 = RGB(135,68,52) = HTML 874434
\usepackage{xcolor}
\usepackage{tikz}
% xcolor 与 l3color 是两套色彩注册表，需各注册同名 sealbrown：
%   xcolor 版 → 供 TikZ/pgf 的 fill= 读取；
%   L3 版     → 供 cls 版框 \color_stroke:n 读取。
\definecolor{sealbrown}{HTML}{874434}
\ExplSyntaxOn
\color_set:nnn { sealbrown } { HTML } { 874434 }
\ExplSyntaxOff

\zhvtset{
  grid color       = sealbrown,
  main font        = GenRyuMin2TW-SB,
  main font path   = C:/Users/Administrator/texmf/fonts/opentype/public/GenRyuMin2TW/,
jiazhu font      = GenRyuMin2TW-M,
  jiazhu font path = C:/Users/Administrator/texmf/fonts/opentype/public/GenRyuMin2TW/,
  % 页码框容纳 4 字：正常页「九十七」+前置空 1 格=4 行恰好；
  % 100+ 页「一百零一」共 5 行超出，末字「一/二/三」被框裁掉（2026-10-04 实测），故放宽到 5
  shukou page number chars = 5,
  % 书口书名框高从默认 6 字减到 4 字，把省出的 2 字高拨给 chapter 框，
  % 使 6 字书名+卷名放进 9 字高的 chapter 框后不再压线（2026-10-05）
  shukou title chars = 4,
}

\title{史记}[五帝本纪]
\author{汉·司马迁}
\maker{闲叟漫记·重製}

% ★ 修正书口卷名位置：chapter 框起始位置改回 cls 默认值，避免全局影响。
%   只在每个卷首页通过 \juan 命令里的空命令微调该页位置（安全，不会跨页污染）。
\ExplSyntaxOn
\AtBeginDocument{
  \dim_set:Nn \l__zhvt_shukou_chapter_title_start_dim
    { \l__zhvt_shukou_title_height_dim + \l__zhvt_font_size_dim * 2 }
  \dim_set:Nn \l__zhvt_shukou_page_number_start_dim
    { \l__zhvt_shukou_chapter_title_start_dim + \l__zhvt_shukou_chapter_title_height_dim }
}
\ExplSyntaxOff

% ★ 述赞样式：题头“索隐述赞”四字用黑体（simhei，与正文宋体/仿宋区分）+ 圆角棕框反白；
%   正文韵文按最初要求用仿宋（simfang），整体缩小字号（21pt）以便每四句单独成列、居列中。
%   ⚠ 竖排关键：自定义字体家族必须带 RawFeature={vertical:+vert}（仅 +vert）。
%      simhei 若加 +vhal 会整体加载失败回退成仿宋；simfang 同理只用 +vert。
\newCJKfontfamily{\shuzanfont}[Path=C:/Windows/Fonts/,RawFeature={vertical:+vert}]{simhei.ttf}
\newCJKfontfamily{\shuzanbody}[Path=C:/Windows/Fonts/,RawFeature={vertical:+vert}]{simfang.ttf}
% 圆角方框题头（2026-11-08 三次修正 · 定稿：字不得越框、不得有重影、墨迹真居中）：
%   ★ 弃用「隐形棕字节点 + 偏移 fill + 白字」三步法：text=sealbrown 的棕字层并非真隐形，
%     露出框外的部分会在白底上显出淡影，视觉上就是「反白字超出底纹」。
%   ★ 改用「定尺寸框 + 单层白字」：框 132pt × 38pt（绝对坐标），圆角 4pt。
%   ★★ 最大陷阱：pymupdf 的 span bbox 是 **advance 字身盒**，而 simhei 竖排替换字形
%      (vert) 的**实际墨迹相对字身盒整体左偏约 18pt**。只看 bbox 会误判「字在框内」，
%      实际笔画已被框左边切断。验收必须做像素级墨迹测量（渲染后统计框内白像素）。
%   ★ 校准结论（像素实测）：白字节点置于 TikZ (0pt, 8.2pt) —— x 方向**不需补偿**
%      （墨迹水平本就居中，早先的 -19.6/-9.6 是拿字身盒当墨迹造成的过补偿）；
%      y 方向 +8.2pt 抵消竖排墨迹垂直偏移。实测墨迹心与框心 dx=0.0 dy=0.0，
%      四边留白 L/R≈21pt、T/B≈8.5pt，且 21pt > 圆角 4pt，绝不切字越框。
%   页面 /Rotate 90 后，横长框自然成为竖条、四字竖排一列。
\newcommand{\shuzanhead}[1]{%
  \begin{tikzpicture}[baseline={(0,2)}]
    \fill[sealbrown, rounded corners=6pt]
      (-76pt,-15pt) rectangle (76pt,15pt);
    \node[text=white, inner sep=0pt, anchor=center] at (0pt,8.2pt)
      {\shuzanfont\fontsize{23}{23}\selectfont #1};
  \end{tikzpicture}%
}

\begin{document}

% ★★ 版心（书口）排布覆写：把书名放到鱼尾之下（2026-10-04）
%
%   cls 默认版心自上而下：书名（title 静态框，在鱼尾之上）→ 鱼尾 → 篇名（chapter 动态框）→ 页码
%   本模板改为：          鱼尾（位置不动）→ 书名 → 篇名 → 页码分割线 → 页码
%
%   手法：把 cls 放在 title 静态框里的书名，改放进 chapter 动态框（鱼尾之下）；
%         title 静态框留空。书名取 \l__zhvt_title_str，即 \title{...} 的第一参数，
%         所以换书时只改 \title，不必在此重复写书名。
%
%   ⚠ 三条必须遵守的写法（2026-10-04 实测）：
%     ① 汉字必须用 \hbox{...} 包住。flowfram 会把框内容装进一个宽 14.4mm（= 一个 line_sep
%        = 书口栏宽）的纵向 minipage 并按该宽度折行，而整页靠 PDF /Rotate 90 显示、
%        框内文字本身不变向 —— 不包 \hbox 时「围炉夜话」4 字会被逐字折成 4 行，
%        在版面上就成了横排 4 字并横向溢出版框。
%     ② 竖排下**一个 \hbox 就是一整列**。书名与篇名要同处一列，必须放进同一个 \hbox；
%        分成两个 \hbox 时，后一个会被当成新的一列、甩到隔壁栏去
%        （实测「序」跑到书口栏外侧 14.4mm 的正文列上）。中间的 \quad 在框内是横向空白，
%        显示出来正好是书名与篇名之间的纵向间隔。
%     ③ ★ 内容盒要写成 **与框等宽的定宽盒** `\hbox to \l__zhvt_shukou_chapter_title_height_dim{…\hfil}`。
%        实测：chapter 框是把内容盒沿纵向**居中**摆放的，内容宽度一变，书名就跟着漂：
%        自然宽度的 \hbox（书名 38.4mm）落在 x=157.2mm，加篇名后（57.6mm）落在 x=152.4mm，
%        目录页（无篇名）与正文页的书名就差半格。做成等宽盒后，无论有没有篇名，
%        内容固定从框顶起排，实测两页均落在 x=142.8mm，不再漂。
\ExplSyntaxOn
% 把书口 chapter 框（书名+篇名）起始位置下移 1 字高，避免 4 字卷名顶到鱼尾
% ⚠ cls 的 \__zhvt_calculate_dimensions: 在 \AtBeginDocument 里执行，会重算该 dim；
%   必须把覆写挂到 begin document 之后，否则会被二次覆盖（2026-10-05 实测）。
% (旧方案已移到导言区，直接修改 chapter 框起始位置)
% 卷标题命令 #1=书口篇名(鱼尾) #2=正文卷名（含卷序，卷名内空格用 ~）
%
% ★★ 目录页两项修订（2026-10-05，愚大人指出）★★
%   ① 页码起点统一对齐：
%      原状 —— 目录项写成 \par\hspace{4em}卷名\hspace{4em}页码，页码起点 = 卷名实宽 + 4em。
%      卷名实宽随字数变：「卷一~梁惠王上」6 字 +~、「卷七~离娄上」5 字 +~，
%      差 1 个字位 → 卷七~卷十 4 条的页码整体比其余 10 条高 27.21pt（一个字），看着上下不齐。
%      改法 —— 卷名补全角空格（宽 = 1 字 = 27.215pt）补足到 **7 字符位**再写进目录，
%      14 条页码起点遂全部落在同一条竖线上（起点对齐；末字随页码字数自然参差）。
%      ⚠ 只对目录里的副本 \l_tmpc_tl 补，正文卷名仍用 \l_tmpa_tl 原样 —— 正文像素不受影响。
%   ② 卷名整体上移半字位：前导 \hspace 由 4em 改 3.5em。
%      实测原状目录块在版框内偏下（天头留白 139.6pt / 地脚留白 110.7pt），
%      上移 0.5em = 13.6pt 后上下基本均分，视觉居中。
\NewDocumentCommand { \juan } { mm } {
	\clearpage
	\tl_set:Nn \l__zhvt_chapter_title_tl { #1 }
	\tl_set:Nn \l_tmpa_tl { #2 }
	\tl_set_eq:NN \l_tmpc_tl \l_tmpa_tl
	\int_step_inline:nn { 7 - \str_count:n { #2 } }
	  { \tl_put_right:Nn \l_tmpc_tl { 　 } }
	\tl_set:Nx \l_tmpb_tl { \zhnum { page } }
	\tl_put_left:Nn \l_tmpb_tl { \hspace { 4em } }
	\tl_put_left:Nn \l_tmpb_tl { \l_tmpc_tl }
	\tl_put_left:Nn \l_tmpb_tl { \par \hspace { 3.5em } }
	\addtocontents { toc } { \l_tmpb_tl }
	\skip_horizontal:n { \l__zhvt_font_size_dim * 4 }
	\l_tmpa_tl
	\par
	\thispagestyle { empty }
}
\cs_set_protected:Npn \__zhvt_set_shukou_content: {
	% ★ 鱼尾上方留白（title 静态框）嵌入闲叟钤印：
	%   整页 /Rotate 90，故此处把图用 \rotatebox{90} 转正为竖章；
	%   尺寸卡在版心条幅内（书内约 7mm 宽 × 33.8mm 高，留余量不压鱼尾）。
	\setstaticcontents * { title } {
		\vbox_to_ht:nn { \l__zhvt_baseline_skip_dim } {
			\vskip 14.2pt
			\hbox_to_wd:nn { \l__zhvt_shukou_title_height_dim } {
				\hfil
				\rotatebox [ origin = c ] { 90 } { \includegraphics [ width = 7mm , height = 33.8mm ] { yshs.png } }
				\hfil
			}
		}
	}
% ★ 版心（书口）统一改法（2026-11-16 愚大人定）：18pt＋框内居中＋书口文字同版框色
%   ① 原 27.2pt 时 7 字长篇名（如「廉颇蔺相如列传」）末字溢出 chapter 框、压住页码；
%      改 \fontsize{18pt}{18pt}，前后 \hfil 令字串在框内居中（不随篇名字数漂移）。
%   ② \vbox{\vskip 8.2pt} 定量下移，抵消缩字号后栏宽方向 8.2pt 上偏（像素实测 y 中心 543.3 = 原版）。
%      ⚠ \raisebox / \vspace* 在 flowfram 框内无效；换字号须同步换补偿：16pt→10.1、19pt→7.4。
%   ③ chapter 与 page~number 均置 \color{sealbrown}（HTML 874434），与版框／界行同色；正文仍黑。
%   ⚠ 页码字位、钤印、鱼尾、界行、正文一律未动；五部合订本同此口径。
	\setdynamiccontents * { chapter }
	{ \color{sealbrown} \vbox { \vskip 8.2pt \hbox to \l__zhvt_shukou_chapter_title_height_dim
		{ \hfil \fontsize{18pt}{18pt}\selectfont \l__zhvt_title_str \quad \l__zhvt_chapter_title_tl \hfil } } }
	\setdynamiccontents * { page~number } { \color{sealbrown} \quad \zhnum { page } }
}
\__zhvt_set_shukou_content:
\ExplSyntaxOff

\maketitle

% 单篇，无多卷目录，不渲染目录页




% 方案①：仅改 A 类错讹，头部与版式不变

\juan{五帝本纪}{卷一~五帝本纪第一}
\hspace*{2em}黄帝者。少典之子。姓公孙。名曰轩辕。生而神灵。弱而能言。幼而徇齐。长而敦敏。成而聪明。\\
\hspace*{2em}轩辕之时。神农氏世衰。诸侯相侵伐。暴虐百姓。而神农氏弗能征。于是轩辕乃习用干戈。以征不享。诸侯咸来宾从。而蚩尤最为暴。莫能伐。炎帝欲侵陵诸侯。诸侯咸归轩辕。轩辕乃修德振兵。治五气。蓺五种。抚万民。度四方。教熊罴貔貅貙虎。以与炎帝战于阪泉之野。三战。然后得其志。蚩尤作乱。不用帝命。于是黄帝乃征师诸侯。与蚩尤战于涿鹿之野。遂禽杀蚩尤。而诸侯咸尊轩辕为天子。代神农氏。是为黄帝。天下有不顺者。黄帝从而征之。平者去之。披山通道。未尝宁居。\\
\hspace*{2em}东至于海。登丸山。及岱宗。西至于空桐。登鸡头。南至于江。登熊湘。北逐荤粥。合符釜山。而邑于涿鹿之阿。迁徙往来无常处。以师兵为营卫。官名皆以云命。为云师。置左右大监。监于万国。万国和。而鬼神。山川。封禅与为多焉。获宝鼎。迎日推策。举风后。力牧。常先。大鸿以治民。顺天地之纪。幽明之占。死生之说。存亡之难。时播百谷草木。淳化鸟兽虫蛾。旁罗日月星辰。水波土石金玉。劳勤心力耳目。节用水火材物。有土德之瑞。故号黄帝。\\
\hspace*{2em}黄帝二十五子。其得姓者十四人。\\
\hspace*{2em}黄帝居轩辕之丘。而娶于西陵之女。是为嫘祖。嫘祖为黄帝正妃。生二子。其后皆有天下。其一曰玄嚣。是为青阳。青阳降居江水。其二曰昌意。降居若水。昌意娶蜀山氏女。曰昌仆。生高阳。高阳有圣德焉。黄帝崩。葬桥山。其孙昌意之子高阳立。是为帝颛顼也。\\
\hspace*{2em}帝颛顼高阳者。黄帝之孙而昌意之子也。静渊以有谋。疏通而知事。养材以任地。载时以象天。依鬼神以制义。治气以教化。絜诚以祭祀。北至于幽陵。南至于交趾。西至于流沙。东至于蟠木。动静之物。大小之神。日月所照。莫不砥属。\\
\hspace*{2em}帝颛顼生子曰穷蝉。颛顼崩。而玄嚣之孙高辛立。是为帝喾。\\
\hspace*{2em}帝喾高辛者。黄帝之曾孙也。高辛父曰蟜极。蟜极父曰玄嚣。玄嚣父曰黄帝。自玄嚣与蟜极皆不得在位。至高辛即帝位。高辛于颛顼为族子。\\
\hspace*{2em}高辛生而神灵。自言其名。普施利物。不于其身。聪以知远。明以察微。顺天之义。知民之急。仁而威。惠而信。修身而天下服。取地之财而节用之。抚教万民而利诲之。历日月而迎送之。明鬼神而敬事之。其色郁郁。其德嶷嶷。其动也时。其服也士。帝喾溉执中而遍天下。日月所照。风雨所至。莫不从服。\\
\hspace*{2em}帝喾娶陈锋氏女。生放勋。娶娵訾氏女。生挚。帝喾崩。而挚代立。帝挚立。不善崩。而弟放勋立。是为帝尧。\\
\hspace*{2em}帝尧者。放勋。其仁如天。其知如神。就之如日。望之如云。富而不骄。贵而不舒。黄收纯衣。彤车乘白马。能明驯德。以亲九族。九族既睦。平章百姓。百姓昭明。合和万国。\\
\hspace*{2em}乃命羲。和。敬顺昊天。数法日月星辰。敬授民时。分命羲仲。居郁夷。曰旸谷。敬道日出。便程东作。日中。星鸟。以殷中春。其民析。鸟兽字微。申命羲叔。居南交。便程南讹。敬致。日永。星火。以正中夏。其民因。鸟兽希革。申命和仲。居西土。曰昧谷。敬道日入。便程西成。夜中。星虚。以正中秋。其民夷易。鸟兽毛毨。申命和叔。居北方。曰幽都。便在伏物。日短。星昴。以正中冬。其民燠。鸟兽氄毛。岁三百六十六日。以闰月正四时。信饬百官。众功皆兴。\\
\hspace*{2em}尧曰。﹁谁可顺此事。﹂放齐曰。﹁嗣子丹朱开明。﹂尧曰。﹁吁。顽凶。不用。﹂尧又曰。﹁谁可者。﹂驩兜曰。﹁共工旁聚布功。可用。﹂尧曰。﹁共工善言。其用僻。似恭漫天。不可。﹂尧又曰。﹁嗟。四岳。汤汤洪水滔天。浩浩怀山襄陵。下民其忧。有能使治者。﹂皆曰鲧可。尧曰。﹁鲧负命毁族。不可。﹂岳曰。﹁异哉。试不可用而已。﹂尧于是听岳用鲧。九岁。功用不成。\\
\hspace*{2em}尧曰。﹁嗟。四岳。朕在位七十载。汝能庸命。践朕位。﹂岳应曰。﹁鄙德忝帝位。﹂尧曰。﹁悉举贵戚及疏远隐匿者。﹂众皆言于尧曰。﹁有矜在民间。曰虞舜。﹂尧曰。﹁然。朕闻之。其何如。﹂岳曰。﹁盲者子。父顽。母嚚。弟傲。能和以孝。烝烝治。不至奸。﹂尧曰。﹁吾其试哉。﹂于是尧妻之二女。观其德于二女。舜饬下二女于妫汭。如妇礼。尧善之。乃使舜慎和五典。五典能从。乃遍入百官。百官时序。宾于四门。四门穆穆。诸侯远方宾客皆敬。尧使舜入山林川泽。暴风雷雨。舜行不迷。尧以为圣。召舜曰。﹁女谋事至而言可绩。三年矣。女登帝位。﹂舜让于德不怿。正月上日。舜受终于文祖。文祖者。尧大祖也。\\
\hspace*{2em}于是帝尧老。命舜摄行天子之政。以观天命。舜乃在璿玑玉衡。以齐七政。遂类于上帝。禋于六宗。望于山川。辩于群神。揖五瑞。择吉月日。见四岳诸牧。班瑞。岁二月。东巡狩。至于岱宗。祡。望秩于山川。遂见东方君长。合时月正日。同律度量衡。修五礼五玉三帛二生一死为挚。如五器。卒乃复。五月。南巡狩。八月。西巡狩。十一月。北巡狩。皆如初。归。至于祖祢庙。用特牛礼。五岁一巡狩。群后四朝。遍告以言。明试以功。车服以庸。肇十有二州。决川。象以典刑。流宥五刑。鞭作官刑。扑作教刑。金作赎刑。眚灾过。赦。怙终贼。刑。钦哉。钦哉。惟刑之静哉。\\
\hspace*{2em}驩兜进言共工。尧曰不可而试之工师。共工果淫辟。四岳举鲧治鸿水。尧以为不可。岳彊请试之。试之而无功。故百姓不便。三苗在江淮。荆州数为乱。于是舜归而言于帝。请流共工于幽陵。以变北狄。放驩兜于崇山。以变南蛮。迁三苗于三危。以变西戎。殛鲧于羽山。以变东夷。四罪而天下咸服。\\
\hspace*{2em}尧立七十年得舜。二十年而老。令舜摄行天子之政。荐之于天。尧辟位凡二十八年而崩。百姓悲哀。如丧父母。三年。四方莫举乐。以思尧。尧知子丹朱之不肖。不足授天下。于是乃权授舜。授舜。则天下得其利而丹朱病。授丹朱。则天下病而丹朱得其利。尧曰﹁终不以天下之病而利一人﹂。而卒授舜以天下。尧崩。三年之丧毕。舜让辟丹朱于南河之南。诸侯朝觐者不之丹朱而之舜。狱讼者不之丹朱而之舜。讴歌者不讴歌丹朱而讴歌舜。舜曰﹁天也﹂。夫而后之中国践天子位焉。是为帝舜。\\
\hspace*{2em}虞舜者。名曰重华。重华父曰瞽叟。瞽叟父曰桥牛。桥牛父曰句望。句望父曰敬康。敬康父曰穷蝉。穷蝉父曰帝颛顼。颛顼父曰昌意。以至舜七世矣。自从穷蝉以至帝舜。皆微为庶人。\\
\hspace*{2em}舜父瞽叟盲。而舜母死。瞽叟更娶妻而生象。象傲。瞽叟爱后妻子。常欲杀舜。舜避逃。及有小过。则受罪。顺事父及后母与弟。日以笃谨。匪有解。\\
\hspace*{2em}舜。冀州之人也。舜耕历山。渔雷泽。陶河滨。作什器于寿丘。就时于负夏。舜父瞽叟顽。母嚚。弟象傲。皆欲杀舜。舜顺适不失子道。兄弟孝慈。欲杀。不可得。即求。尝在侧。\\
\hspace*{2em}舜年二十以孝闻。三十而帝尧问可用者。四岳咸荐虞舜。曰可。于是尧乃以二女妻舜以观其内。使九男与处以观其外。舜居妫汭。内行弥谨。尧二女不敢以贵骄事舜亲戚。甚有妇道。尧九男皆益笃。舜耕历山。历山之人皆让畔。渔雷泽。雷泽上人皆让居。陶河滨。河滨器皆不苦窳。一年而所居成聚。二年成邑。三年成都。尧乃赐舜絺衣。与琴。为筑仓廪。予牛羊。瞽叟尚复欲杀之。使舜上涂廪。瞽叟从下纵火焚廪。舜乃以两笠自扞而下。去。得不死。后瞽叟又使舜穿井。舜穿井为匿空旁出。舜既入深。瞽叟与象共下土实井。舜从匿空出。去。瞽叟。象喜。以舜为已死。象曰﹁本谋者象。﹂象与其父母分。于是曰。﹁舜妻尧二女。与琴。象取之。牛羊仓廪予父母。﹂象乃止舜宫居。鼓其琴。舜往见之。象鄂不怿。曰。﹁我思舜正郁陶。﹂舜曰。﹁然。尔其庶矣。﹂舜复事瞽叟爱弟弥谨。于是尧乃试舜五典百官。皆治。\\
\hspace*{2em}昔高阳氏有才子八人。世得其利。谓之﹁八恺﹂。高辛氏有才子八人。世谓之﹁八元﹂。此十六族者。世济其美。不陨其名。至于尧。尧未能举。舜举八恺。使主后土。以揆百事。莫不时序。举八元。使布五教于四方。父义。母慈。兄友。弟恭。子孝。内平外成。\\
\hspace*{2em}昔帝鸿氏有不才子。掩义隐贼。好行凶慝。天下谓之浑沌。少暤氏有不才子。毁信恶忠。崇饰恶言。天下谓之穷奇。颛顼氏有不才子。不可教训。不知话言。天下谓之檮杌。此三族世忧之。至于尧。尧未能去。缙云氏有不才子。贪于饮食。冒于货贿。天下谓之饕餮。天下恶之。比之三凶。舜宾于四门。乃流四凶族。迁于四裔。以御螭魅。于是四门辟。言毋凶人也。\\
\hspace*{2em}舜入于大麓。烈风雷雨不迷。尧乃知舜之足授天下。尧老。使舜摄行天子政。巡狩。舜得举用事二十年。而尧使摄政。摄政八年而尧崩。三年丧毕。让丹朱。天下归舜。而禹。皋陶。契。后稷。伯夷。夔。龙。倕。益。彭祖自尧时而皆举用。未有分职。于是舜乃至于文祖。谋于四岳。辟四门。明通四方耳目。命十二牧论帝德。行厚德。远佞人。则蛮夷率服。舜谓四岳曰。﹁有能奋庸美尧之事者。使居官相事。﹂皆曰。﹁伯禹为司空。可美帝功。﹂舜曰。﹁嗟。然。禹。汝平水土。维是勉哉。﹂禹拜稽首。让于稷。契与皋陶。舜曰。﹁然。往矣。﹂舜曰。﹁弃。黎民始饥。汝后稷播时百谷。﹂舜曰。﹁契。百姓不亲。五品不驯。汝为司徒。而敬敷五教。在宽。﹂舜曰。﹁皋陶。蛮夷猾夏。寇贼奸轨。汝作士。五刑有服。五服三就。五流有度。五度三居。维明能信。﹂舜曰。﹁谁能驯予工。﹂皆曰垂可。于是以垂为共工。舜曰。﹁谁能驯予上下草木鸟兽。﹂皆曰益可。于是以益为朕虞。益拜稽首。让于诸臣朱虎。熊罴。舜曰。﹁往矣。汝谐。﹂遂以朱虎。熊罴为佐。舜曰。﹁嗟。四岳。有能典朕三礼。﹂皆曰伯夷可。舜曰。﹁嗟。伯夷。以汝为秩宗。夙夜维敬。直哉维静絜。﹂伯夷让夔。龙。舜曰。﹁然。以夔为典乐。教稺子。直而温。宽而栗。刚而毋虐。简而毋傲。诗言意。歌长言。声依永。律和声。八音能谐。毋相夺伦。神人以和。﹂夔曰。﹁于。予击石拊石。百兽率舞。﹂舜曰。﹁龙。朕畏忌谗说殄伪。振惊朕众。命汝为纳言。夙夜出入朕命。惟信。﹂舜曰。﹁嗟。女二十有二人。敬哉。惟时相天事。﹂三岁一考功。三考绌陟。远近众功咸兴。分北三苗。\\
\hspace*{2em}此二十二人咸成厥功。皋陶为大理。平。民各伏得其实。伯夷主礼。上下咸让。垂主工师。百工致功。益主虞。山泽辟。弃主稷。百谷时茂。契主司徒。百姓亲和。龙主宾客。远人至。十二牧行而九州莫敢辟违。唯禹之功为大。披九山。通九泽。决九河。定九州。各以其职来贡。不失厥宜。方五千里。至于荒服。南抚交趾。北发。西戎。析枝。渠廋。氐。羌。北山戎。发。息慎。东长。鸟夷。四海之内咸戴帝舜之功。于是禹乃兴九招之乐。致异物。凤皇来翔。天下明德皆自虞帝始。\\
\hspace*{2em}舜年二十以孝闻。年三十尧举之。年五十摄行天子事。年五十八尧崩。年六十一代尧践帝位。践帝位三十九年。南巡狩。崩于苍梧之野。葬于江南九疑。是为零陵。舜之践帝位。载天子旗。往朝父瞽叟。夔夔唯谨。如子道。封弟象为诸侯。舜子商均亦不肖。舜乃豫荐禹于天。十七年而崩。三年丧毕。禹亦乃让舜子。如舜让尧子。诸侯归之。然后禹践天子位。尧子丹朱。舜子商均。皆有疆土。以奉先祀。服其服。礼乐如之。以客见天子。天子弗臣。示不敢专也。\\
\hspace*{2em}自黄帝至舜。禹。皆同姓而异其国号。以章明德。故黄帝为有熊。帝颛顼为高阳。帝喾为高辛。帝尧为陶唐。帝舜为有虞。帝禹为夏后而别氏。姓姒氏。契为商。姓子氏。弃为周。姓姬氏。\\
\hspace*{2em}太史公曰。学者多称五帝。尚矣。然尚书独载尧以来。而百家言黄帝。其文不雅驯。荐绅先生难言之。孔子所传宰予问五帝德及帝系姓。儒者或不传。余尝西至空桐。北过涿鹿。东渐于海。南浮江淮矣。至长老皆各往往称黄帝。尧。舜之处。风教固殊焉。总之不离古文者近是。予观春秋。国语。其发明五帝德。帝系姓章矣。顾弟弗深考。其所表见皆不虚。书缺有间矣。其轶乃时时见于他说。非好学深思。心知其意。固难为浅见寡闻道也。余并论次。择其言尤雅者。故著为本纪书首。\\
% ===== 述赞（索隐述赞）=====
% 题头「索隐述赞」四字：黑体 + 圆角棕框反白，已居中（见 \shuzanhead）。
% 正文韵文：仿宋、21pt；每四句为一列，各列在列内居中对齐。
%   ★ 定宽盒换列法（2026-11-08 实测三种方案后唯一全对者）：
%     每组包成 \hbox to \hsize{\hfil 组文 \hfil} —— 盒宽=列宽，盒内 \hfil 使组文
%     在列内严格居中；盒间用 \linebreak 提供断点（★\hspace*{30em} 强制换列会吞掉
%     下一列首项，\linebreak 不会；盒间裸接 % 无断点会溢出被裁）。
%     题头后也用 \linebreak（同理不能用 \hspace*{30em}）。
%   - 首字「曰/高/爰/就/明」全部存活（实测 col1–col5 组心均为 456.5pt，严丝合缝对齐）。
%   - ★ 列距必须 = 41pt（\fontsize{24}{41}）：全书界行网格为主文列相位 121+41k、
%     每页 21 列；旧 21/25 的 25pt 列距让述赞文字落在两线之间、直接压住界行线。
%     改 41pt 后每列正好落在格子中心，字不再压线；字号 24pt 与正文 27.2pt 略有层次。
%   - 述赞整体另起一页（\clearpage）：实测不加分页时题头+第1组挤在正文末页末 2 列、
%     后 4 组被拆到下一页；加分页后 6 列同页完整呈现。
\clearpage
\shuzanhead{索隐述赞曰}\linebreak
{\shuzanbody\fontsize{24}{41}\selectfont
\hbox to \hsize{\hfil 帝出少典。居于轩丘。既代炎历。遂禽蚩尤。\hfil}\linebreak
\hbox to \hsize{\hfil 高阳嗣位。静深有谋。小大远近。莫不怀柔。\hfil}\linebreak
\hbox to \hsize{\hfil 爰洎帝喾。列圣同休。帝挚之弟。其名放勋。\hfil}\linebreak
\hbox to \hsize{\hfil 就之如日。望之如云。郁夷东作。昧谷西曛。\hfil}\linebreak
\hbox to \hsize{\hfil 明扬侧陋。玄德升闻。能让天下。贤哉二君。\hfil}%
\par}
\end{document}
